\section{Categorical preliminaries}\label{sec:preliminaries}

\subsection{Partial maps and restriction idempotents}

Composition is written $gf=g\circ f$, so $f$ is applied first.
A \emph{restriction category}, in the sense of Cockett and
Lack~\cite[Section~2.1.1]{CockettLack}, is a category equipped with an assignment
that sends every morphism
$f:X\rightharpoonup Y$ to a morphism
$\bar f:X\rightharpoonup X$ satisfying
\[
 f\bar f=f,\qquad
 \bar f\bar g=\bar g\bar f,\qquad
 \overline{g\bar f}=\bar g\bar f,\qquad
 \bar h f=f\overline{hf}.
\]
In the middle two equations $g:X\rightharpoonup Z$; in the last,
$h:Y\rightharpoonup Z$.
The morphisms $\bar f$ are idempotent and are called
\emph{restriction idempotents}, or \emph{domains} in
recursion-category terminology \cite[Section~5]{Heller}.
The morphism $f$ is \emph{total} when $\bar f=1_X$.
For a restriction category $\mathcal C$, let
$\operatorname{Tot}(\mathcal C)$ be its \emph{category of total morphisms}:
the subcategory with all objects of $\mathcal C$ and its
total morphisms.
A \emph{restriction functor} $F:\mathcal C\to\mathcal D$ between
restriction categories satisfies $F(\bar f)=\overline{F(f)}$ for every
morphism $f$. A \emph{restriction equivalence} consists of restriction
functors in both directions whose composites are naturally isomorphic
to the respective identity functors.

Restriction idempotents on $X$ are ordered by $a\le b$ when $ba=a$.
Their meet is composition, their greatest element is $1_X$, and
$a\vee b$ denotes their least upper bound when it exists.
For parallel partial morphisms, the restriction order is
$f\le g$ when $f=g\bar f$.
Parallel morphisms $f,g:X\rightharpoonup Y$ are \emph{compatible} when
$f\bar g=g\bar f$. Their join, when it exists, is their least upper
bound in the restriction order.
A \emph{restriction zero} $0_{X,Y}:X\rightharpoonup Y$ is
absorbing under composition and satisfies
$\overline{0_{X,Y}}=0_{X,X}$. We abbreviate $0_{X,X}$ to $0_X$.
Restriction idempotents $a,b$ on $X$ are disjoint when $ab=0_X$.
They are complementary when also $a\vee b=1_X$.
A disjoint pair $a_0,a_1$ is \emph{inseparable} if no complementary
pair $c,c'$ satisfies $a_0\le c$ and $a_1\le c'$.

For $f:X\rightharpoonup Y$ and a restriction idempotent $a$ on $Y$,
the \emph{inverse image} of $a$ is the restriction idempotent
$f^*a=\overline{af}$ on $X$.
A restriction idempotent $r$ on $Y$ \emph{receives} $f$ when $rf=f$
\cite[Section 6]{Heller}.
The \emph{range} $\ran(f)$ is the least restriction idempotent receiving $f$;
we require the agreement law
\[
 gf=hf\quad\Longrightarrow\quad
 g\ran(f)=h\ran(f)
 \qquad(g,h:Y\rightharpoonup Z).
\]
For a partial function $f$, the restriction idempotent
$\bar f$ and the range $\ran(f)$ are the partial identities on its
domain of definition and image, respectively.

A \emph{restriction monomorphism} $m:A\to X$ is a total morphism
with a partial inverse $r:X\rightharpoonup A$ satisfying
$rm=1_A$ and $mr=\bar r$.
A restriction idempotent $e$ on $X$ \emph{splits} if
$e=mr$ for a total $m:A\to X$ and $r:X\rightharpoonup A$
with $rm=1_A$. The category is \emph{split} when every restriction
idempotent splits.
The \emph{restriction idempotent completion} of a restriction category
$\mathcal C$, denoted by $\operatorname{Split}(\mathcal C)$, is obtained
by splitting the restriction idempotents of $\mathcal C$.
Cockett--Lack denote this category by $K_r(\mathcal C)$
\cite[Section~2.3.3]{CockettLack}.
The objects of $\operatorname{Split}(\mathcal C)$ are pairs $(X,e)$, where $e$ is a
restriction idempotent on $X$. Its morphisms
$f:(X,e)\rightharpoonup(Y,d)$ are morphisms
$f:X\rightharpoonup Y$ satisfying $f=dfe$.
The identity of $(X,e)$ is $e$, and $f$ is total there when
$\bar f=e$. Thus $(X,e)$ is the object through which $e$ splits.

The \emph{Karoubi envelope} $\operatorname{Kar}(\mathcal C)$ has the
same description with arbitrary idempotents $e:X\rightharpoonup X$
in place of restriction idempotents. The restriction operation on
$\operatorname{Kar}(\mathcal C)$ sends
$f:(X,e)\rightharpoonup(Y,d)$ to the morphism represented by $\bar f e$.
This is the construction $K_E(\mathcal C)$ of
\cite[Section~2.3.3, Proposition~2.26]{CockettLack} with $E$ the class
of all idempotents.

For a category $\mathcal E$ with pullbacks, let
$\operatorname{Par}(\mathcal E,\operatorname{Mono})$ be its category
of partial maps, where $\operatorname{Mono}$ denotes all monomorphisms
of $\mathcal E$ \cite[Section~3]{CockettLack}.
Its objects are those of $\mathcal E$. A morphism
$X\rightharpoonup Y$ is an isomorphism class of spans
$X\xleftarrow{m}D\xrightarrow{f}Y$ with $m$ monic, and composition
uses pullback. Its restriction is represented by
$X\xleftarrow{m}D\xrightarrow{m}X$.
We identify its total category with $\mathcal E$.

\subsection{Products, coproducts, and local connectedness}

A \emph{binary restriction product} $X\times Y$ has total
projections $\pi_1:X\times Y\to X$ and $\pi_2:X\times Y\to Y$.
For morphisms $f:Z\rightharpoonup X$ and $g:Z\rightharpoonup Y$,
there is a unique $\langle f,g\rangle:Z\rightharpoonup X\times Y$
with
\[
 \pi_1\langle f,g\rangle=f\bar g,\qquad
 \pi_2\langle f,g\rangle=g\bar f,\qquad
 \overline{\langle f,g\rangle}=\bar f\bar g.
\]
Products of morphisms are obtained by this pairing and the projections.
The total diagonal is
$\Delta_X=\langle1_X,1_X\rangle:X\to X\times X$.
A restriction category has a $P$-category structure if and only if
it has binary restriction products
\cite[Proposition 4.5]{CockettLackIII}.

Binary \emph{restriction coproducts} are categorical coproducts with
total injections $i_0:X\to X+Y$ and $i_1:Y\to X+Y$.
For morphisms $f:X\rightharpoonup Z$ and $g:Y\rightharpoonup Z$,
write $[f,g]:X+Y\rightharpoonup Z$ for their copairing.
The \emph{codiagonal} is
$\nabla_X=[1_X,1_X]:X+X\to X$.
In a category with ranges and restriction coproducts,
$a\vee b=\ran[a,b]$ for restriction idempotents on the same object.

A range functor in the sense of Cockett, Guo, and Hofstra
\cite[p.~424]{CGHRangeI} preserves restrictions and ranges. Following
\cite[p.~7]{Lengyel2006}, we also require a \emph{range functor}
$F:\mathcal C\to\mathcal D$ to preserve
restriction zeros and binary restriction coproducts. In particular, the canonical morphism
$\chi_{X,Y}=[F i_0,F i_1]:FX+FY\to F(X+Y)$
is an isomorphism, and
$F[f,g]\chi_{X,Y}=[Ff,Fg]$ for $f:X\rightharpoonup Z$ and
$g:Y\rightharpoonup Z$. A range functor preserves binary
joins of restriction idempotents by the preceding formula, and hence
complementary pairs.
It need not preserve restriction products. Let $\Pfn$ denote the
category of sets and partial functions. For example, the functor
$F:\Pfn\to\Pfn$ given by $F(X)=X\times\{0,1\}$ and
$F(f)=f\times1_{\{0,1\}}$ preserves restrictions, images, zeros,
and coproducts. Its canonical morphism from $F(X\times Y)$ to
$FX\times FY$ sends $((x,y),a)$ to $((x,a),(y,a))$, and is
not surjective when $X,Y$ are nonempty.

A restriction idempotent $a$ on $X$ is \emph{connected} if, for
restriction idempotents $b,c$ on $X$, the equations $a=b\vee c$ and
$bc=0$ imply $b=0$ or $c=0$. The category is \emph{locally connected}
when its connected restriction idempotents are jointly epimorphic: for
$f,g:X\rightharpoonup Y$, the equations $fa=ga$ for every connected
restriction idempotent $a$ on $X$ imply $f=g$.
A \emph{partial monomorphism} $p:X\rightharpoonup Y$ satisfies
$pu=pv\Rightarrow\bar p u=\bar p v$ for all
$u,v:W\rightharpoonup X$.
A \emph{section} of $p$ is a morphism $s:Y\rightharpoonup X$
with $ps=\bar s$ and $psp=p$.
\emph{Weak choice} means that every partial monomorphism has such
a section. For partial functions, this is the existence of the
partial inverse of an injection on its domain.

\subsection{Recursion categories}

A restriction category with binary restriction products and zeros is
\emph{prodominical} if $f\times0=0$ for every morphism $f$ and every
zero morphism of the required type. It is \emph{isotypical} if any
two of its objects are isomorphic.

A \emph{Turing morphism} on $U$ is a morphism
$\tau:U\times U\rightharpoonup U$ such that, for every object $X$
and every $f:X\times U\rightharpoonup U$, there is a total
parameter morphism $r:X\to U$ satisfying $f=\tau(r\times1_U)$.
A \emph{recursion category} is a prodominical isotypical
category with binary restriction products and a Turing morphism,
as in \cite[Section 9]{Heller}.

In a restriction category $\mathcal C$, a \emph{restriction terminal object}
$1$ has a unique total morphism $!_X:X\to1$ for each object $X$,
satisfying $!_Yf=!_X\bar f$ for every $f:X\rightharpoonup Y$.
A restriction category is \emph{cartesian} if it has binary
restriction products and a restriction terminal object.
A \emph{Turing category} is a cartesian restriction category with an
object $U$, called a \emph{Turing object}, such that every object $X$
admits a total morphism $m_X:X\to U$ and a morphism
$r_X:U\rightharpoonup X$ with $r_Xm_X=1_X$, and there is a
Turing morphism on $U$. This is the characterization of
Cockett--Hofstra~\cite[Theorem 3.4]{CockettHofstra}.


